\documentclass{situationist-preprint}

\title{Dialectical Effectiveness Versus Soundness}
\msid{EPI-02-c}
\msfamily{EPI-02}
\msdiscipline{Epistemology}
\mstype{Consequence}
\mskeywords{Moore's proof; transmission failure; dialectic; begging the question; scepticism}

\begin{document}
\maketitle

\begin{abstract}
Wright's charge that Moore's proof fails to transmit warrant is often taken to show that the proof is defective. We distinguish three properties of an argument: soundness, the transmission of knowledge from premises to conclusion, and dialectical effectiveness against a given interlocutor. From the conclusion of EPI-02-a, that Moore's proof is sound and that Moore knows its premises and its conclusion by means of it, we derive that the proof has the first two properties and lacks the third against anyone who doubts that he is not deceived. The phenomena that motivate the transmission-failure diagnosis are phenomena of the third kind. Refuting scepticism, in the sense of EPI-02, does not require persuading the sceptic.
\end{abstract}

\section{Three properties}

\citet{Moore1939} held up his hands and argued: here is one hand, here is another; so there are at least two external things. Many readers have thought that the proof, though valid with true premises, is useless. \citet{Wright2002} explained the impression by a principle of transmission failure: if one's warrant for a premise depends on an antecedent entitlement to accept the conclusion, or something that entails it, then the argument cannot be one's means of acquiring warrant for the conclusion. Moore's warrant for \enquote{here is a hand} depends, on Wright's view, on an entitlement to take it that he is not dreaming or envatted, so the proof cannot give him warrant for anything that entails that he is not.

Three properties of an argument must be kept apart.

\begin{definition}\label{d:three}
An argument is \emph{sound} if it is valid and its premises are true. It \emph{transmits knowledge} for a subject if the subject knows its premises and comes to know its conclusion by competent deduction from them. It is \emph{dialectically effective} against an interlocutor if the interlocutor, given what she already accepts and the evidence she is willing to treat as such, is rationally required to accept its premises, and hence its conclusion.
\end{definition}

Soundness and transmission concern the relation between the argument, the world and the arguer. Dialectical effectiveness concerns a relation between the argument and someone else's state of opinion.

\section{The consequences}

\begin{proposition}\label{p:has}
Moore's proof is sound and, in ordinary contexts, transmits knowledge for Moore, both to its official conclusion and to the conclusion that he is not envatted.
\end{proposition}

\begin{proof}
Soundness is Proposition~3.1 of EPI-02-a. By Theorem~3.2 of EPI-02-a, Moore knows in ordinary contexts that he has hands, and knows that there are external things and that he is not envatted on bases of competent deduction from that knowledge.
\end{proof}

\begin{proposition}\label{p:lacks}
Moore's proof is not dialectically effective against an interlocutor who suspends judgment on whether Moore is envatted.
\end{proposition}

\begin{proof}
Such an interlocutor treats the hypothesis that Moore is envatted as open. On that hypothesis Moore would have the same perceptual basis for \enquote{here is a hand} as he actually has (Assumption~2.7 of EPI-02-a). So nothing she accepts and nothing she counts as evidence favours the premise over its denial, and she is not rationally required to accept it. Since she is not required to accept the premise, the argument does not require her to accept its conclusion.
\end{proof}

Propositions~\ref{p:has} and~\ref{p:lacks} are consistent, and their conjunction is what Wright's examples exhibit. The examples show that the proof cannot move someone who does not already accept that Moore is not envatted. They do not show that Moore lacks knowledge of the premise, or that his knowledge of the premise depends on antecedent knowledge of the conclusion. On the account of EPI-02-a, it does not: perceptual justification is immediate (EPI-19), and the further condition knowledge requires, safety, is a condition on the world rather than a further belief. There is therefore nothing for the warrant for \enquote{I am not envatted} to be antecedent to. As \citet{Pryor2004} argued, the proof's defect is dialectical: it is question-begging against a certain audience, and that is a property of the proof relative to the audience, not of the proof.

\begin{corollary}\label{c:refute}
To refute scepticism, in the sense established in EPI-02-a, is to show that the sceptical argument has a false premise, and this can be done by an argument that is not dialectically effective against the sceptic.
\end{corollary}

\section{Objection}

\begin{objection}
An anti-sceptical argument that cannot move the sceptic is no answer to her. Philosophy is a dialectical enterprise, and an argument that convinces only those who already agree has no philosophical use.
\end{objection}

\begin{reply}
Three uses remain. First, the argument settles a question of fact for those who are not sceptics: whether the sceptical argument is sound. Second, it tells the reflective non-sceptic that her ordinary knowledge claims are not undermined by her inability to answer the sceptic, which is the source of the unease the sceptic produces. Third, it diagnoses the sceptic's position: her suspension of judgment on whether she is envatted is what makes the proof ineffective against her, and that suspension is not required by her evidence, which is what EPI-02-a shows. No argument with premises about the external world could be dialectically effective against an interlocutor who suspends judgment on all of them; that this is so is a fact about such interlocutors, and does not count against the arguments.
\end{reply}

\section{Conclusion}

Moore's proof is sound and transmits knowledge; it does not persuade the sceptic. The charge of transmission failure conflates the second property with the third. Once they are separated, the dialectical weakness of the proof is seen to be compatible with the claim of EPI-02-a that scepticism is refuted.

\begin{thebibliography}{9}\small
\bibitem[Moore(1939)]{Moore1939} Moore, G.~E. (1939). Proof of an external world. \emph{Proceedings of the British Academy}, 25, 273--300.
\bibitem[Pryor(2004)]{Pryor2004} Pryor, J. (2004). What's wrong with Moore's argument? \emph{Philosophical Issues}, 14, 349--378.
\bibitem[Wright(2002)]{Wright2002} Wright, C. (2002). (Anti-)sceptics simple and subtle: G.~E. Moore and John McDowell. \emph{Philosophy and Phenomenological Research}, 65(2), 330--348.
\end{thebibliography}

\end{document}
